% 第4章 4.1.3：D delta_0 的测试反例。phi_n(x)=chi(x) sin(nx)，chi 为光滑截断（[-0.3,0.3]上为1，支撑[-0.9,0.9]）。
% n=3（绿虚线）与 n=10（蓝实线）：sup|phi_n|<=1，支撑固定，而原点切线斜率 phi_n'(0)=n（短切线段）。数据由同目录 generate.py 计算。
\begin{tikzpicture}[font=\figurefont\footnotesize]
\begin{axis}[width=72mm,height=40mm,axis lines=left,
  x axis line style={draw=FigureStroke,line width=.6pt},y axis line style={draw=FigureStroke,line width=.6pt},
  tick style={draw=FigureStroke},tick label style={font=\figurefont\scriptsize,text=FigureInk},
  xmin=-1.1,xmax=1.1,ymin=-1.25,ymax=1.25,xtick={-0.9,0,0.9},ytick={-1,0,1},
  xlabel={$x$},xlabel style={at={(axis description cs:1,0.5)},anchor=west,font=\figurefont\scriptsize},
  extra y ticks={0},extra y tick style={grid=major,major grid style={draw=FigureGrid,line width=.4pt},tick label style={opacity=0}},
  legend style={draw=none,fill=none,font=\figurefont\scriptsize,at={(1.02,1)},anchor=north west,cells={anchor=west}}]
\def\dat{figures/02-abstract-mathematics/03-function-spaces/tikz/c4-ddelta-test/tests.dat}
\addplot[draw=FigureBlue,line width=1.1pt] table[x=x,y=n10] {\dat};
\addplot[draw=FigureGreen,line width=1.1pt,dashed] table[x=x,y=n3] {\dat};
\addplot[draw=FigureInk,line width=.7pt,domain=-0.09:0.09] {10*x};
\addplot[draw=FigureInk,line width=.7pt,domain=-0.3:0.3] {3*x};
\legend{$\varphi_{10}$,$\varphi_{3}$,原点切线}
\end{axis}
\end{tikzpicture}
